\chapter{Reduced dataset analysis}

In this section, we analyze the behavior of the Logistic Regression model when training data is severely limited. By keeping only 1 out of every 50 samples from the original training set, we drastically reduce the number of observations available for parameter estimation. This setup allows us to highlight the vital role of explicit regularization as a tool to prevent overfitting in low-data regimes.

\section{Experimental results}

The model was trained on the reduced dataset ($DTR_{reduced}, LTR_{reduced}$) using the exact same sweep of 13 $\lambda$ values as the standard analysis. The validation was still performed on the full validation set to ensure a reliable and comparable estimate of the generalization error for the target application $\pi_T = 0.1$.

The numerical results are summarized in Table \ref{tab:red_lr_results} and the visual trends are plotted in Figure \ref{fig:lr_reduced_plot}.

\begin{table}[h]
\centering
\caption{Reduced dataset Logistic Regression results ($\pi_T = 0.1$)}
\label{tab:red_lr_results}
\begin{tabular}{@{}ccc@{}}
\toprule
\textbf{Regularization Coefficient ($\lambda$)} & \textbf{minDCF} & \textbf{actDCF} \\ \midrule
$1.0000 \times 10^{-4}$ & 0.4466 & 0.9780 \\
$3.1623 \times 10^{-4}$ & 0.4487 & 0.9780 \\
$1.0000 \times 10^{-3}$ & 0.4487 & 0.7466 \\
$3.1623 \times 10^{-3}$ & 0.4431 & \textbf{0.4652} \\
$1.0000 \times 10^{-2}$ & 0.4407 & \textbf{0.4652} \\
$3.1623 \times 10^{-2}$ & 0.4407 & \textbf{0.4652} \\
$1.0000 \times 10^{-1}$ & 0.3988 & 0.7164 \\
$3.1623 \times 10^{-1}$ & 0.3874 & 0.9856 \\
$1.0000 \times 10^{0}$  & 0.3874 & 1.0000 \\
$3.1623 \times 10^{0}$  & 0.3874 & 1.0000 \\
$1.0000 \times 10^{1}$  & \textbf{0.3783} & 1.0000 \\
$3.1623 \times 10^{1}$  & 0.3807 & 1.0000 \\
$1.0000 \times 10^{2}$  & 0.3920 & 1.0000 \\
\bottomrule
\end{tabular}
\end{table}

\begin{figure}[h]
    \centering
    \includegraphics[width=0.85\textwidth]{LR_Standard_Reduced_DCF_vs_Lambda.png}
    \caption{Actual DCF and Minimum DCF as a function of $\lambda$ for the Standard Logistic Regression model trained on the reduced dataset (1/50 samples).}
    \label{fig:lr_reduced_plot}
\end{figure}

The optimal discriminative capability is reached at a highly regularized $\lambda = 10^{1}$, yielding a minDCF of \textbf{0.3783}. Conversely, the best calibration occurs much earlier, between $\lambda = 3.16 \times 10^{-3}$ and $3.16 \times 10^{-2}$, where the actDCF drops to a minimum of \textbf{0.4652}.

\section{Analysis and Considerations}

By comparing Figure \ref{fig:lr_reduced_plot} with the standard plot from the previous chapter, we can observe entirely different model dynamics dictated by the scarcity of data. We can break this behavior down into three distinct regularization regimes:

\begin{itemize}
    \item \textbf{Low regularization ($\lambda \le 10^{-3}$): Overfitting.} 
    Unlike the full dataset model where minDCF was flat, here a low $\lambda$ produces terrible results for both metrics (minDCF $\approx 0.448$, actDCF $\approx 0.978$). With very few training samples and virtually no penalty on the weights, the model completely memorizes the noise in the tiny training set. The resulting weights become extremely large and highly tailored to those specific points, destroying the model's ability to generalize to the evaluation set. The model produces wildly overconfident, miscalibrated scores, ruining both its ranking ability and its absolute thresholding.

    \item \textbf{Moderate regularization ($\lambda \approx 10^{-2}$): Calibration recovery.} 
    As we increase $\lambda$, we penalize extreme weight values. This forces the model to ignore spurious correlations (noise) and rely only on the most robust, generalizable features. This variance reduction significantly improves the reliability of the scores, bringing the empirical risk closer to the true risk. This results in the best-calibrated scores and the lowest actDCF ($0.4652$).

    \item \textbf{High regularization ($\lambda \ge 10^{-1}$): Underfitting and calibration loss.} 
    As $\lambda$ increases further, the penalty term heavily dominates the cross-entropy loss. The weights $w$ are aggressively pushed toward zero, causing the model to become too rigid to capture the underlying data distribution (underfitting). Just as in the standard model, these "shrunken" weights compress the scores $s(x) = w^Tx + b$ toward a constant value dictated entirely by the bias. The scores lose their probabilistic meaning, forcing the system to predict the negative class for all samples. Consequently, the actDCF hits the maximum ceiling of $1.0000$.
\end{itemize}

\subsection*{Why does minDCF continue to improve at high $\lambda$?}
An interesting phenomenon in the reduced dataset is that the minDCF actually reaches its absolute best value ($0.3783$) at $\lambda = 10^{1}$, well into the regime where actDCF is completely broken ($1.0000$). 

This happens because minDCF only cares about the \textit{relative sorting} of the samples. In a severely data-starved environment, heavily squashing the weights limits the model to acting essentially like a mean-distance classifier. While the absolute values of these heavily regularized scores are useless for direct thresholding (hence actDCF = 1), their relative ranking is actually more stable and robust than the erratic, overfitted ranking produced by a low $\lambda$.

\section{Theoretical Conclusion}
In conclusion, when training data is abundant, the model is naturally regularized by the sheer volume of information, rendering the L2 penalty mostly irrelevant for discriminative power. However, when data is scarce, explicit regularization via $\lambda$ is strictly mandatory. It is the only mechanism available to control the bias-variance tradeoff, preventing catastrophic overfitting and allowing the model to recover a minDCF ($0.3783$) that is surprisingly close to the full-dataset baseline ($0.3611$).